% !TEX program = xelatex
% demo 讲义的源文件（用 xelatex 编译后即为 <实验标题>/<实验标题>.pdf）
% 真实使用时：把课程下发的讲义 PDF 直接放进 <实验标题>/ 并改成同名即可，无需本文件。
\documentclass[11pt]{ctexart}
\usepackage[margin=2.4cm]{geometry}
\usepackage{amsmath,booktabs}
\title{单摆法测重力加速度}
\author{近代物理实验（演示讲义）}
\date{2026 年春季学期}
\begin{document}
\maketitle

\section*{一、实验目的}
\begin{enumerate}
  \item 掌握用单摆测量重力加速度的原理与方法；
  \item 学会用最小二乘直线处理数据并评估线性程度；
  \item 分析计时误差、摆长测量误差对结果的影响。
\end{enumerate}

\section*{二、实验原理}
摆长 $L$ 远大于摆球半径、摆角 $\theta_0<5^\circ$ 时，单摆作简谐振动，周期
\begin{equation}
  T=2\pi\sqrt{\frac{L}{g}},
\end{equation}
其中 $g$ 为当地重力加速度。若测量 $n$ 个周期的累计时间 $t_n$，则 $T=t_n/n$。由此
\begin{equation}
  g=\frac{4\pi^{2}L}{T^{2}} .
\end{equation}
为减小单点计算的偶然误差，将上式改写为 $T^{2}=\dfrac{4\pi^{2}}{g}L$，以 $T^{2}$ 对 $L$ 作最小二乘直线 $T^{2}=kL+b$，则 $g=4\pi^{2}/k$。

\section*{三、仪器与器材}
单摆装置（可调摆线 + 摆球）、米尺（分度值 1 mm）、游标卡尺、光电门计时器或秒表。

\section*{四、实验内容}
\begin{enumerate}
  \item 测量摆球直径 $d$，摆长按 $L=l+d/2$ 计算（$l$ 为悬点到球顶的线长）；
  \item 使 $L$ 依次取 $0.500$、$0.600$、$0.700$、$0.800$、$0.900\ \mathrm{m}$；
  \item 每个摆长下测 50 个周期的累计时间 $t_{50}$，各测一次并记录；
  \item 以 $T^{2}$ 对 $L$ 作最小二乘直线，求斜率 $k$、截距 $b$、相关系数 $R^2$，并计算 $g=4\pi^{2}/k$；
  \item 与当地重力加速度参考值 $9.79\ \mathrm{m/s^2}$ 比较，给出相对偏差。
\end{enumerate}

\section*{五、数据处理要求}
\begin{itemize}
  \item 原始记录须完整保留（摆长、$t_{50}$），不得涂改；
  \item 派生量（$T$、$T^{2}$、逐次 $g$）由程序计算，报告须附计算脚本；
  \item 报告需给出拟合直线图、残差量级与误差来源分析。
\end{itemize}

\end{document}
